\begin{proof}[Proof of \cref{obj:lem:observable-factorization}]
Fix \(z\in\mathbb C\).  The definitions give
\[
Z_n=T-\bar g_n(X)=\eta+\{g_0(X)-\bar g_n(X)\}=\eta+D_n(X).
\]
Moreover, \cref{obj:def:plm-model,obj:def:non-gaussian-class} give the pointwise clipping bound \(|\bar g_n(X)|\le C_g\) and the treatment-regression range bound \(|g_0(X)|\le C_g\) almost surely, hence
\[
 |D_n(X)|=|g_0(X)-\bar g_n(X)|
 \le |g_0(X)|+|\bar g_n(X)|
 \le 2C_g
 \qquad P\text{-a.s.}
\]
Together with the Luxemburg treatment-tail condition for \(\eta\), this bound gives every real exponential moment of \(Z_n\), and hence the exponential integrability used below. % lean: learnedResidual_integrable_exp
Since \(D_n(X)\) is a function of \(X\) and \(\eta\) is independent of \(X\) under \cref{obj:def:non-gaussian-class},
\[
\begin{aligned}
F_n(z)
&=E_P\exp\{z(\eta+D_n(X))\}  \\
&=E_P\bigl[\exp(z\eta)\exp\{zD_n(X)\}\bigr] \\
&=E_P\exp(z\eta)\;E_P\exp\{zD_n(X)\}
=M(z)H_n(z).
\end{aligned}
\]
% lean: residualMGF_eq_treatmentMGF_mul_nuisanceMGF

For the second identity, first note the pointwise residual decomposition
\[
Y=\theta_0 Z_n+b_n(X)+\xi,
\]
because \cref{obj:def:outcome-contamination} gives \(b_n(X)=q_0(X)-\theta_0D_n(X)\), while \(Z_n=\eta+D_n(X)\) and \(\xi=Y-q_0(X)-\theta_0\eta\).  The preceding exponential-moment envelope for \(Z_n\) also gives differentiation under the expectation:
\[
F_n'(z)
=
E_P\!\left[Z_n e^{zZ_n}\right].
\]
% lean: hasDerivAt_complexMGF

Thus
\[
\begin{aligned}
G_n(z)
&=E_P\!\left[Y e^{zZ_n}\right] \\
&=\theta_0 E_P\!\left[Z_ne^{zZ_n}\right]
  +E_P\!\left[b_n(X)e^{zZ_n}\right]
  +E_P\!\left[\xi e^{zZ_n}\right].
\end{aligned}
\]
The contamination weight is bounded almost surely by
\[
\begin{aligned}
 |b_n(X)|
 &=|q_0(X)-\theta_0D_n(X)| \\
 &\le |q_0(X)|+|\theta_0|\,|D_n(X)| \\
 &\le C_q+C_{\theta}(2C_g),
\end{aligned}
\qquad P\text{-a.s.}
\]
where the last line uses the outcome-regression and coefficient range bounds in \cref{obj:def:non-gaussian-class} together with the displayed bound on \(D_n\).  This boundedness and the same exponential envelope for \(Z_n\) give integrability of the contamination term. % lean: contamination_weighted_exp_integrable

The outcome-noise term is integrable as well.  The \(\xi\)-tail condition in \cref{obj:def:non-gaussian-class} gives \(E_P[|\xi|^2]<\infty\), and the learned-residual envelope gives
\[
E_P\!\left[\left|e^{zZ_n}\right|^2\right]
=E_P\!\left[e^{2\operatorname{Re}(z)Z_n}\right]<\infty.
\]
By Cauchy--Schwarz, \(E_P[|\xi e^{zZ_n}|]<\infty\). % lean: outcomeNoise_weighted_exp_integrable
Since \(Z_n=T-\bar g_n(X)\), the real and imaginary parts of \(e^{zZ_n}\) are measurable with respect to \((X,T)\).  Applying \(E_P[\xi\mid X,T]=0\) from \cref{obj:def:non-gaussian-class} separately to these integrable real products yields
\[
E_P\!\left[\xi e^{zZ_n}\right]=0.
\]
% lean: outcomeNoise_weighted_exp_eq_zero

Finally, independence of \(\eta\) and \(X\) factors the contamination term:
\[
\begin{aligned}
E_P\!\left[b_n(X)e^{zZ_n}\right]
&=E_P\!\left[b_n(X)e^{zD_n(X)}e^{z\eta}\right] \\
&=E_Pe^{z\eta}\;E_P\!\left[b_n(X)e^{zD_n(X)}\right]
=M(z)B_n(z).
\end{aligned}
\]
% lean: contamination_weighted_factorization

Substituting the last three displays and the derivative identity gives
\[
G_n(z)=\theta_0F_n'(z)+M(z)B_n(z).
\]
Since \(z\) was arbitrary, both identities hold for every \(z\in\mathbb C\). % lean: observable_factorization
\end{proof}
